\section{Related Work}
\subsection{Chinese Food Recognition}
Food-101 \cite{bossard2014food101} remains the standard Western benchmark, and
DeepFood \cite{liu2016deepfood} established that ImageNet pre-training plus
fine-tuning beats hand-crafted pipelines for dietary assessment. Chinese menus
are covered by separate resources: ChineseFoodNet \cite{chen2017chinesefoodnet}
collected 185\,628 images over 208 categories, and CNFOOD-241 \cite{cnfood241},
the dataset used here, provides 241 categories with a fine-grained study of
category-level discrimination \cite{xiong2026cnfood}. Surveys of food computing
\cite{zhang2023foodsurvey} list cross-domain generalization, portion estimation
and multi-dish scenes as the open problems, which is exactly the setting of a
canteen tray; our contribution is not a new dataset but a controlled measurement
on an existing one, with the deployment constraints made explicit.

\subsection{Transfer Learning and Pre-Training Scale}
Features learned on ImageNet \cite{deng2009imagenet} transfer well to
fine-grained targets, but the benefit depends on source-target similarity and on
how many layers are adapted \cite{yosinski2014transfer}, and better source models
do not always transfer better \cite{kornblith2019better}. The usual lever in a
course-scale project is adaptation depth; we keep that lever fixed and vary the
source domain instead, replacing ImageNet features with features learned on 241
Chinese dishes. The comparison is only meaningful because the recipe, the head
capacity and the schedule are identical between the two runs.

\subsection{Open-Set Recognition}
Open-set recognition asks a classifier to reject inputs from classes it was not
trained on \cite{scheirer2013openset}, and OpenMax \cite{bendale2016openset}
introduced the activation-based formulation that most later work builds on. Our
gate is deliberately simpler and closer to a deployment artefact than to a
benchmark method: a logistic probe on pooled features, a veto from a wider head,
and a trained catch-all class. We evaluate it by acceptance rate on five sets
rather than by a single open-set score, because the quantity a user experiences
is whether the system answers, not the area under a curve.

\subsection{Calibration and Explainability}
Modern networks are over-confident \cite{guo2017calibration}, which matters when
a system reports an energy range derived from its confidence, so we report
expected calibration error and fit one temperature per branch on validation.
Robustness to common corruptions \cite{hendrycks2019robustness} motivates the
proxy domain of the baseline branch. Grad-CAM \cite{selvaraju2017gradcam}
provides the localization maps we use to check that the network attends to food
rather than to the tray.
